% figures/w12-sector-chords.svg — 한 단위의 실현 가능 영역은 **원의 부채꼴**이고,
% QP 는 직선만 받으므로 원호를 **내접 현**으로 바꾼다.
%
%   bash _tools/tikz2svg.sh figures/src/w12-sector-chords.tex figures/w12-sector-chords.svg
%
% 사용자 지시 (2026-11-04): "각 n rpm값의 saturation과 서보 각도의 saturation도
% 고려해야 하고...."  — 벽이 둘이고, 둘이 서로 다르게 생겼다는 것이 §12-12 다.
%
% 두 칸이 §12-12 의 두 줄이다.
%  (왼쪽)  서보 벽은 직선 둘, 추력 캡은 원호 하나. 합치면 부채꼴이다.
%          원호는 2차식이므로 QP 의 제약으로 쓸 수 없다.
%  (오른쪽) 원호를 각을 M 등분한 현 M 개로 바꾼다. 이웃한 현은 반지름 R 의 원
%          **위**에서 만나므로 다각형은 내접하고, 중점에서만 0.966 R 이다.
%
% 좌표는 눈대중이 아니다 (`CLAUDE.md` §4-6-1) — `_tools/w12_sector_chords_geom.m`
% 이 찍은 값을 그대로 옮겼다. alpha_max = 45 도, M = 3, R = 2.30 cm 에서
%   꼭짓점  -45: (1.62635, -1.62635)   -15: (2.22163, -0.59528)
%           +15: (2.22163, +0.59528)   +45: (1.62635, +1.62635)
%   현의 중점 -30: (1.92399, -1.11081)   0: (2.22163, 0)   +30: (1.92399, +1.11081)
%   R cos15 = 0.9659 R,  중점에서 3.41 % 보수적
%   확대경   틈 0.07837 cm, 원본 원의 중심 (-30 도, 2.26081), 반지름 0.0775 cm,
%            x8 로 키우면 0.627 cm — 반지름 0.62 cm 의 확대경 안에 들어간다
% 식은 `W12_alloc` 의 `cone_cap` 과 같은 것이다. 바꿀 때 둘을 함께 고친다.
%
% 앞 그림들에서 배운 것 셋을 여기서도 지킨다 (`standing-orders.md` §15-41).
%  (1) 칸 둘의 중심 간격(7.6 cm)은 제목이 서로 닿지 않을 만큼 넓게. 7.0 cm 에서는
%      제목 둘이 겹쳤다 — 간격을 넓히고 제목을 줄여서 고쳤다.
%  (2) 제목 노드는 축 이름표보다 5 mm 이상 위에. 축을 2.40 에서 끊고 이름표를
%      2.22 에 두어 부제(3.10)와 떨어뜨렸다.
%  (3) `cap` 은 TikZ 내장 키라 스타일 이름으로 쓸 수 없다 — `capl` 로 둔다.
%
% 틈이 R 의 3.4 % 라 본문 폭에서 2 px 다. **보이지 않는 것을 숫자로만 주장하지
% 않는다** (`CLAUDE.md` §4-6) — -30 도 현의 중점을 x8 로 확대해 그 틈을 그린다.
\documentclass[border=5pt]{standalone}
% gnc-style.tex — 이 볼트의 TikZ 그림이 공유하는 스타일.
%
% 저널 그림 규약을 스타일 하나로 굳혀 둔다. 그림마다 선 굵기나 화살촉을
% 다시 정하지 않는다 — 그것이 그림이 제각각으로 보이는 원인이다.
%
% 쓰는 법:  % gnc-style.tex — 이 볼트의 TikZ 그림이 공유하는 스타일.
%
% 저널 그림 규약을 스타일 하나로 굳혀 둔다. 그림마다 선 굵기나 화살촉을
% 다시 정하지 않는다 — 그것이 그림이 제각각으로 보이는 원인이다.
%
% 쓰는 법:  % gnc-style.tex — 이 볼트의 TikZ 그림이 공유하는 스타일.
%
% 저널 그림 규약을 스타일 하나로 굳혀 둔다. 그림마다 선 굵기나 화살촉을
% 다시 정하지 않는다 — 그것이 그림이 제각각으로 보이는 원인이다.
%
% 쓰는 법:  \input{gnc-style}  (figures/src/ 안에서)

\usepackage{tikz}
\usepackage{amsmath}
\usepackage{xcolor}
\usetikzlibrary{arrows.meta, positioning, calc, fit, backgrounds, decorations.pathmorphing}

% ---- 색 --------------------------------------------------------------------
% 색은 강조 하나에만 쓴다. 나머지는 검정.
\definecolor{ink}{HTML}{1A1A1A}
\definecolor{accent}{HTML}{7C3AED}
\definecolor{warn}{HTML}{B91C1C}
\definecolor{soft}{HTML}{666666}

% 기하 그림(좌표계·유도)에서만 쓰는 넷. 블록선도는 위의 흑백 + accent 로 충분하다.
\definecolor{pathc}{HTML}{0D9488}   % 계획 경로 · 항적
\definecolor{errc}{HTML}{C2410C}    % 오차 (cross-track 등)
\definecolor{alng}{HTML}{B45309}    % along-track
\definecolor{flow}{HTML}{0369A1}    % 조류 · 유속

% ---- 선과 화살촉 -----------------------------------------------------------
\tikzset{
  % 신호선. 화살촉은 작고 뾰족한 것 하나뿐이다.
  sig/.style   = {ink, line width=0.5pt, -{Stealth[length=4pt, width=3pt]}},
  wire/.style  = {ink, line width=0.5pt},
  % 강조 경로 — 파선으로 구분한다. 흑백 인쇄에서도 살아야 하기 때문이다.
  asig/.style  = {accent, line width=0.5pt, densely dashed,
                  -{Stealth[length=4pt, width=3pt]}},
  awire/.style = {accent, line width=0.5pt, densely dashed},
  % 블록. 흰 바탕, 직각, 검은 테두리.
  blk/.style   = {draw=ink, line width=0.55pt, fill=white,
                  minimum width=17mm, minimum height=8mm, inner sep=2pt, font=\small},
  % 합산점
  sum/.style   = {draw=ink, line width=0.55pt, fill=white, circle,
                  inner sep=0pt, minimum size=4.6mm, font=\scriptsize},
  asum/.style  = {draw=accent, line width=0.55pt, fill=white, circle,
                  inner sep=0pt, minimum size=4.6mm, font=\scriptsize},
  % 분기점
  dot/.style   = {ink, fill=ink, circle, inner sep=0pt, minimum size=1.5mm},
  % 신호 이름 — 선 위에 작게
  sname/.style = {font=\small, inner sep=1.5pt},
  % 그림 안의 짧은 설명. 그림 안의 글은 최소로 둔다
  note/.style  = {font=\scriptsize, soft, inner sep=1.5pt},
  anote/.style = {font=\scriptsize, accent, inner sep=1.5pt},
  % 축
  axis/.style  = {ink, line width=0.5pt},
  tick/.style  = {font=\scriptsize, soft},
}

% 캡션 한 줄 — 그림 아래 가는 선과 함께
\newcommand{\gnccaption}[2]{%
  \node[anchor=north west, text width=#1, align=left, font=\scriptsize, ink]
        at (current bounding box.south west) {\vspace{2pt}#2};}
  (figures/src/ 안에서)

\usepackage{tikz}
\usepackage{amsmath}
\usepackage{xcolor}
\usetikzlibrary{arrows.meta, positioning, calc, fit, backgrounds, decorations.pathmorphing}

% ---- 색 --------------------------------------------------------------------
% 색은 강조 하나에만 쓴다. 나머지는 검정.
\definecolor{ink}{HTML}{1A1A1A}
\definecolor{accent}{HTML}{7C3AED}
\definecolor{warn}{HTML}{B91C1C}
\definecolor{soft}{HTML}{666666}

% 기하 그림(좌표계·유도)에서만 쓰는 넷. 블록선도는 위의 흑백 + accent 로 충분하다.
\definecolor{pathc}{HTML}{0D9488}   % 계획 경로 · 항적
\definecolor{errc}{HTML}{C2410C}    % 오차 (cross-track 등)
\definecolor{alng}{HTML}{B45309}    % along-track
\definecolor{flow}{HTML}{0369A1}    % 조류 · 유속

% ---- 선과 화살촉 -----------------------------------------------------------
\tikzset{
  % 신호선. 화살촉은 작고 뾰족한 것 하나뿐이다.
  sig/.style   = {ink, line width=0.5pt, -{Stealth[length=4pt, width=3pt]}},
  wire/.style  = {ink, line width=0.5pt},
  % 강조 경로 — 파선으로 구분한다. 흑백 인쇄에서도 살아야 하기 때문이다.
  asig/.style  = {accent, line width=0.5pt, densely dashed,
                  -{Stealth[length=4pt, width=3pt]}},
  awire/.style = {accent, line width=0.5pt, densely dashed},
  % 블록. 흰 바탕, 직각, 검은 테두리.
  blk/.style   = {draw=ink, line width=0.55pt, fill=white,
                  minimum width=17mm, minimum height=8mm, inner sep=2pt, font=\small},
  % 합산점
  sum/.style   = {draw=ink, line width=0.55pt, fill=white, circle,
                  inner sep=0pt, minimum size=4.6mm, font=\scriptsize},
  asum/.style  = {draw=accent, line width=0.55pt, fill=white, circle,
                  inner sep=0pt, minimum size=4.6mm, font=\scriptsize},
  % 분기점
  dot/.style   = {ink, fill=ink, circle, inner sep=0pt, minimum size=1.5mm},
  % 신호 이름 — 선 위에 작게
  sname/.style = {font=\small, inner sep=1.5pt},
  % 그림 안의 짧은 설명. 그림 안의 글은 최소로 둔다
  note/.style  = {font=\scriptsize, soft, inner sep=1.5pt},
  anote/.style = {font=\scriptsize, accent, inner sep=1.5pt},
  % 축
  axis/.style  = {ink, line width=0.5pt},
  tick/.style  = {font=\scriptsize, soft},
}

% 캡션 한 줄 — 그림 아래 가는 선과 함께
\newcommand{\gnccaption}[2]{%
  \node[anchor=north west, text width=#1, align=left, font=\scriptsize, ink]
        at (current bounding box.south west) {\vspace{2pt}#2};}
  (figures/src/ 안에서)

\usepackage{tikz}
\usepackage{amsmath}
\usepackage{xcolor}
\usetikzlibrary{arrows.meta, positioning, calc, fit, backgrounds, decorations.pathmorphing}

% ---- 색 --------------------------------------------------------------------
% 색은 강조 하나에만 쓴다. 나머지는 검정.
\definecolor{ink}{HTML}{1A1A1A}
\definecolor{accent}{HTML}{7C3AED}
\definecolor{warn}{HTML}{B91C1C}
\definecolor{soft}{HTML}{666666}

% 기하 그림(좌표계·유도)에서만 쓰는 넷. 블록선도는 위의 흑백 + accent 로 충분하다.
\definecolor{pathc}{HTML}{0D9488}   % 계획 경로 · 항적
\definecolor{errc}{HTML}{C2410C}    % 오차 (cross-track 등)
\definecolor{alng}{HTML}{B45309}    % along-track
\definecolor{flow}{HTML}{0369A1}    % 조류 · 유속

% ---- 선과 화살촉 -----------------------------------------------------------
\tikzset{
  % 신호선. 화살촉은 작고 뾰족한 것 하나뿐이다.
  sig/.style   = {ink, line width=0.5pt, -{Stealth[length=4pt, width=3pt]}},
  wire/.style  = {ink, line width=0.5pt},
  % 강조 경로 — 파선으로 구분한다. 흑백 인쇄에서도 살아야 하기 때문이다.
  asig/.style  = {accent, line width=0.5pt, densely dashed,
                  -{Stealth[length=4pt, width=3pt]}},
  awire/.style = {accent, line width=0.5pt, densely dashed},
  % 블록. 흰 바탕, 직각, 검은 테두리.
  blk/.style   = {draw=ink, line width=0.55pt, fill=white,
                  minimum width=17mm, minimum height=8mm, inner sep=2pt, font=\small},
  % 합산점
  sum/.style   = {draw=ink, line width=0.55pt, fill=white, circle,
                  inner sep=0pt, minimum size=4.6mm, font=\scriptsize},
  asum/.style  = {draw=accent, line width=0.55pt, fill=white, circle,
                  inner sep=0pt, minimum size=4.6mm, font=\scriptsize},
  % 분기점
  dot/.style   = {ink, fill=ink, circle, inner sep=0pt, minimum size=1.5mm},
  % 신호 이름 — 선 위에 작게
  sname/.style = {font=\small, inner sep=1.5pt},
  % 그림 안의 짧은 설명. 그림 안의 글은 최소로 둔다
  note/.style  = {font=\scriptsize, soft, inner sep=1.5pt},
  anote/.style = {font=\scriptsize, accent, inner sep=1.5pt},
  % 축
  axis/.style  = {ink, line width=0.5pt},
  tick/.style  = {font=\scriptsize, soft},
}

% 캡션 한 줄 — 그림 아래 가는 선과 함께
\newcommand{\gnccaption}[2]{%
  \node[anchor=north west, text width=#1, align=left, font=\scriptsize, ink]
        at (current bounding box.south west) {\vspace{2pt}#2};}


\definecolor{knownc}{HTML}{1D4ED8}

\tikzset{
  ax/.style     = {ink, line width=0.5pt, -{Stealth[length=4pt,width=3pt]}},
  wall/.style   = {accent, line width=1.0pt},
  capl/.style   = {warn, line width=1.0pt},
  chord/.style  = {knownc, line width=1.0pt},
  ghost/.style  = {warn, line width=0.7pt, densely dotted},
  lead/.style   = {soft, line width=0.4pt},
  shade/.style  = {fill=accent!12},
  sliver/.style = {fill=warn!45},
  dim/.style    = {ink, line width=0.4pt,
                   {Stealth[length=3pt,width=2.5pt]}-{Stealth[length=3pt,width=2.5pt]}},
  lab/.style    = {font=\scriptsize, align=left,   inner sep=1.5pt},
  clab/.style   = {font=\scriptsize, align=center, inner sep=1.5pt},
  tlab/.style   = {font=\small,      align=center, inner sep=2pt},
  box/.style    = {font=\scriptsize, align=center, inner sep=3pt,
                   draw=knownc, fill=white, rounded corners=1pt, text=knownc},
}

\begin{document}
\begin{tikzpicture}[x=1cm, y=1cm]

% ===== 왼쪽: 부채꼴 =========================================================
\begin{scope}[shift={(0,0)}]
  \fill[shade] (0,0) -- (-45:2.30) arc (-45:45:2.30) -- cycle;

  \draw[ax] (-0.55,0) -- (3.00,0);
  \draw[ax] (0,-2.60) -- (0,2.40);
  \node[clab, anchor=west]       at (3.02,0)     {$x_i'$};
  \node[clab, anchor=south east] at (-0.08,2.22) {$y_i'$};

  % 서보 벽 — 직선 둘 / the servo wall, two straight lines
  \draw[wall] (0,0) -- (45:2.30);
  \draw[wall] (0,0) -- (-45:2.30);
  % 추력 캡 — 원호 하나 / the thrust cap, one arc
  \draw[capl] (-45:2.30) arc (-45:45:2.30);

  % alpha_max 와 R / the angle and the radius
  \draw[ink, line width=0.4pt] (0.95,0) arc (0:45:0.95);
  \node[clab] at (22.5:1.26) {$\alpha_{\max}$};
  \draw[lead] (0,0) -- (-22.5:2.30);
  \node[clab, fill=white, inner sep=1pt] at (-22.5:1.30) {$R$};

  \node[lab, anchor=south west, text=accent] at (0.26,1.82)
    {the servo wall: \textbf{two straight lines}\\
     $\lvert y_i'\rvert \le \tan\alpha_{\max}\, x_i'$};
  \node[lab, anchor=north west, text=warn] at (0.26,-1.86)
    {the thrust cap: \textbf{one arc}\\
     $x_i'^2 + y_i'^2 \le R^2$ --- not linear};

  \node[tlab] at (1.60,3.50) {one unit's feasible set is a \textbf{sector}};
  \node[clab] at (1.60,3.10)
    {the angle is the servo wall, the radius is the shaft-speed cap};
\end{scope}

% ===== 오른쪽: 현으로 바꾼다 ================================================
\begin{scope}[shift={(7.6,0)}]
  % 다각형 / the polygon the QP actually sees
  \fill[shade] (0,0) -- (1.62635,-1.62635) -- (2.22163,-0.59528)
                     -- (2.22163, 0.59528) -- (1.62635, 1.62635) -- cycle;
  % 버린 얇은 조각 셋 / the three slivers that are given up
  \fill[sliver] (-45:2.30) arc (-45:-15:2.30) -- cycle;
  \fill[sliver] (-15:2.30) arc (-15: 15:2.30) -- cycle;
  \fill[sliver] ( 15:2.30) arc ( 15: 45:2.30) -- cycle;

  \draw[ax] (-0.55,0) -- (3.00,0);
  \draw[ax] (0,-2.60) -- (0,2.40);
  \node[clab, anchor=west]       at (3.02,0)     {$x_i'$};
  \node[clab, anchor=south east] at (-0.08,2.22) {$y_i'$};

  % 대신한 원호를 점선으로 남겨 둔다 / the arc it replaced, kept as a dotted line
  \draw[ghost] (-45:2.30) arc (-45:45:2.30);

  \draw[wall] (0,0) -- (45:2.30);
  \draw[wall] (0,0) -- (-45:2.30);
  \draw[chord] (1.62635,-1.62635) -- (2.22163,-0.59528);
  \draw[chord] (2.22163,-0.59528) -- (2.22163, 0.59528);
  \draw[chord] (2.22163, 0.59528) -- (1.62635, 1.62635);

  % 꼭짓점은 원 위에 있다 / the vertices sit on the circle
  \foreach \a in {-45,-15,15,45} { \fill[ink] (\a:2.30) circle (1.6pt); }
  \draw[lead] (2.52,1.14) -- (2.22163,0.59528);
  \node[lab, anchor=west] at (2.54,1.18) {the vertices lie\\\textbf{on} the circle};

  \node[lab, anchor=south west, text=knownc] at (0.26,1.82)
    {the arc becomes $M=3$ \textbf{chords}\\
     $\cos\theta_j\,x_i' + \sin\theta_j\,y_i' \le R\cos(\alpha_{\max}/M)$\\
     $\theta_j = -30^\circ,\ 0^\circ,\ +30^\circ$ \quad
     $R\cos 15^\circ = 0.966\,R$};

  % ---- 확대경: 현과 원호 사이의 틈 / the lens on the chord-to-arc gap ------
  \coordinate (Pc) at (-30:2.26081);
  \coordinate (Lc) at (3.05,-1.95);
  \draw[lead] (Pc) circle (0.0775);
  \draw[lead] ($(Pc)!0.0775cm!(Lc)$) -- ($(Lc)!0.62cm!(Pc)$);
  \begin{scope}[shift={(Lc)}, rotate=-30]
    \begin{scope}
      \clip (0,0) circle (0.62);
      \fill[sliver] (-0.3135,-0.62) rectangle (0.3135,0.62);
      \draw[chord] (-0.3135,-0.62) -- (-0.3135,0.62);
      \draw[ghost] ( 0.3135,-0.62) -- ( 0.3135,0.62);
    \end{scope}
    \draw[lead] (0,0) circle (0.62);
    \draw[dim] (-0.3135,0) -- (0.3135,0);
  \end{scope}
  \node[clab] at (3.05,-2.84) {$\times 8$: chord to arc, $3.4\,\%$ of $R$};

  \node[tlab] at (1.60,3.50) {the QP sees \textbf{chords}, not the arc};
  \node[clab] at (1.60,3.10) {straight lines only, and inscribed in the sector};
\end{scope}

% ===== 아래: 같은 그림이 두 캡 높이를 모두 쓴다 =============================
\node[box] at (5.70,-3.86)
  {In unit axes $x_i'=\sigma_i f_{x,i}$ and $y_i'=\sigma_i f_{y,i}$,
   so one picture serves a unit pushing ahead and one pushing astern.\\
   Only $R$ differs: ahead $R=k_{\mathrm{pos}}n_{\max}^2=119.68$~N,
   astern $R=k_{\mathrm{neg}}n_{\min}^2=66.71$~N, which is $56\,\%$ of it,\\
   so the astern cap is the wall that binds.
   \texttt{W12\_alloc} builds exactly these rows in \texttt{cone\_cap}:
   two for the wall and $M$ for the cap, for each of the four units.};

\end{tikzpicture}
\end{document}
